\section{IBP identities, row operations, and exactness}
\label{app:derivation}
The starting total-derivative identity follows from integration by parts in a region where the integral converges, followed by dimensional analytic continuation~\cite{ChetyrkinTkachov,Weinzierl2022}. Write $P_{\boldsymbol\sigma}=\prod_iD_i^{-\sigma_i}$ and $\partial_k=\partial/\partial\ell_k$. The product rule gives
\begin{align}
 \partial_k\!\cdot(vP_{\boldsymbol\sigma})
 &=(\partial_k\!\cdot v)P_{\boldsymbol\sigma}
 -\sum_i\sigma_i(v\!\cdot\partial_kD_i)
 D_i^{-\sigma_i-1}\prod_{j\ne i}D_j^{-\sigma_j}.
 \label{eq:product-rule}
\end{align}
For $v=\ell_h$, the first factor is $d\delta_{kh}$; for an external momentum it is zero. Completeness of the inverse-propagator coordinates permits us to write $v\!\cdot\partial_kD_i=u_i+\sum_jv_{ij}D_j$, where $u_i$ and $v_{ij}$ are coefficients fixed by the propagators and kinematics, with the invariant and mass terms included in $u_i$. Integrating Eq.~\eqref{eq:product-rule} therefore gives
\begin{equation}
 0=(\partial_k\!\cdot v)I_{\boldsymbol\sigma}
 -\sum_i\sigma_i u_i I_{\boldsymbol\sigma+\boldsymbol e_i}
 -\sum_{i,j}\sigma_i v_{ij}
 I_{\boldsymbol\sigma+\boldsymbol e_i-\boldsymbol e_j}.
 \label{eq:index-shifts-new}
\end{equation}
The unit vector $\boldsymbol e_i$ raises index $i$ by one. This derivation explains how a seed and an operator generate a row of the reduction system.

For the introductory tadpole example, put $D=\ell^2-x+i0$ and $v=\ell$. Direct differentiation gives
\begin{equation}
 \partial_\ell\!\cdot(\ell D^{-a})=dD^{-a}-2a\ell^2D^{-a-1}.
 \label{eq:tadpole-derivative-proof}
\end{equation}
Using $\ell^2=D+x$ with the infinitesimal prescription understood, then integrating, gives $(d-2a)T_a-2axT_{a+1}=0$. Solving at $a=1$ proves Eq.~\eqref{eq:intro-recurrence}. This uses the same analytically continued total-derivative identity as Eq.~\eqref{eq:identity}.

Let $q$ be a source row and $k$ a pivot column with $A_{qk}\ne0$. Solving the identity $\sum_jA_{qj}I_j=0$ for $I_k$ gives $I_k=-\sum_{j\ne k}(A_{qj}/A_{qk})I_j$. Subtracting $A_{rk}/A_{qk}$ times that zero identity from row $r$ yields
\begin{equation}
 A'_{rj}=A_{rj}-\frac{A_{rk}}{A_{qk}}A_{qj},\qquad A'_{rk}=0.
 \label{eq:row-new}
\end{equation}
Every step preserves the represented relation. Repetition until only prescribed master columns remain gives Eq.~\eqref{eq:reduction}. Removing seeds or equations preserves the validity of the retained identities, but can prevent this elimination from finishing. We therefore require reduction to the prescribed basis and exact agreement with the reference coefficients. Generic kinematics avoid singular pivot denominators; special kinematic limits require their own treatment.


For independent components, write the loop momenta as disjoint groups and the integrand as $\prod_b f_b$. In a common convergence domain, repeated integration gives $I=\prod_b\int f_b$. If the chosen integer powers have no such domain in $d$, first vary the propagator powers independently and then continue them to the requested values. Analytic continuation preserves the resulting identity~\cite{Weinzierl2022}. Substituting each exact component reduction therefore gives
\begin{equation}
 I=\prod_b\left(\sum_m r_{bm}M_{bm}\right).
 \label{eq:product-reduction}
\end{equation}
Our product output retains this factored form. It avoids expanding the Cartesian product of master terms and uses the same representation in all timed methods.

\section{The seed regions and equation score}
\label{app:sparse}
For seed $\boldsymbol\sigma$, define excess positive power $D=\sum_i\max(\sigma_i-1,0)$, numerator rank $R=\sum_i\max(-\sigma_i,0)$, and positive sector $P=\{i:\sigma_i>0\}$. Let $h$ be the number of absent physical lines relative to the largest top sector, and let $D_T,R_T$ be the largest target values. The distance from the target numerator rectangles is defined by
\begin{equation}
 q(\boldsymbol\sigma,\mathcal T)=\min_{\boldsymbol a\in\mathcal T}
 \sum_{i\notin P}\max\{0,-\sigma_i-\max(0,-a_i)\}.
 \label{eq:seed-distance}
\end{equation}
Projection onto the interval from zero to each target numerator power leaves precisely this positive excess. Summing the distances and minimizing over targets gives Eq.~\eqref{eq:seed-distance}. We define the region by
\begin{align}
 q&\leq\max\{0,\lfloor\theta_0+\theta_1h+\theta_2D+\theta_6R\rfloor\},\nonumber\\
 R&\leq\max\{0,\lfloor R_T+\theta_3-\theta_4h\rfloor\},\qquad
 D\leq\max\{0,\lfloor D_T+\theta_5\rfloor\}.
 \label{eq:seed-regions}
\end{align}
These inequalities specify a search family rather than a sufficiency theorem. Native outer bounds, retained sector corners, and admissible target and master seeds also apply. The previously studied sparse choices use
\begin{align}
 \boldsymbol\theta_{\rm wide}&=(2.5,0,0,1.25,0,1,-0.25),\nonumber\\
 \boldsymbol\theta_{\rm tight}&=(1,0,0,0,1,1,0),\nonumber\\
 \boldsymbol\theta_{\rm medium}&=(1.5,0,0,0.5,0,1,0).
 \label{eq:seed-values}
\end{align}
The retained native model uses these three sparse choices. The rejected replacement uses the wide and tight choices and a third action that retains all seeds within the tight procedure's native outer bounds. The native procedure retains its subsector truncation and recovery rules. The model selects one complete procedure before generating equations.

The separately tested equation classifier has seventeen features: four family and request dimensions; three seed counts; two differences from target power maxima; three sector descriptors; exact target membership and minimum index distance; the operator index; and the two target maxima. Positive labels mark equations retained by a complete reference reduction. Sampling at most 2,000 records of each class $y$ in batch $c$ gives each sampled record weight $N_{cy}/(n_{cy}N_c)$, where $N_{cy}$ and $n_{cy}$ are the native and sampled class counts and $N_c$ is the full batch count. Multiplying this weight by $n_{cy}$ preserves the class fraction $N_{cy}/N_c$. A leaf with weighted totals $W_0,W_1$ returns $p=W_1/(W_0+W_1)$, and a fixed threshold retains equations. The decision-tree construction follows Ref.~\cite{CART1984}. This empirical score does not certify sufficiency or probability calibration; the native completion check remains necessary.

\section{Measured-cost learning}
\label{app:shared}
Let $i$ label a training request and $a$ an allowed arithmetic procedure. In the latest study, $C_{ia}$ is the measured complete time, including recovery, and $b_i>0$ is the time for the preceding procedure on the same request. Define the relative training target by
\begin{equation}
 y_{ia}=C_{ia}/b_i.
 \label{eq:reconstruction-cost}
\end{equation}
There are five actions, so the latest model predicts the vector $\boldsymbol z_i=(y_{i1},\ldots,y_{i5})$. These quantities are training definitions. The three-loop structure uses the preceding procedure by an explicit guard and is excluded from the 36-request fit. The earlier reconstruction study used six actions and an estimated recovery penalty; its separate objective and outcomes remain in the notes.

The ensemble has 32 extremely randomized regression trees~\cite{Geurts2006ExtraTrees}. At a terminal branch, or leaf, let $S$ be the set of training requests that reach it and let $\mu$ be the predicted value of one output component $z$. We use the loss $\sum_{i\in S}b_i(z_i-\mu)^2$, which defines the baseline-time weighting. Differentiating gives
\begin{equation}
 0=2\sum_{i\in S}b_i(\mu-z_i),\qquad
 \widehat z(S)=\frac{\sum_{i\in S}b_i z_i}{\sum_{i\in S}b_i}.
 \label{eq:reconstruction-leaf}
\end{equation}
Substituting the relative cost from Eq.~\eqref{eq:reconstruction-cost} gives
\begin{equation}
 \widehat y_a(S)=\frac{\sum_{i\in S}C_{ia}}{\sum_{i\in S}b_i},
 \qquad
 \arg\min_a\widehat y_a(S)=\arg\min_a\sum_{i\in S}C_{ia}.
 \label{eq:reconstruction-leaf-choice}
\end{equation}
The second equality follows because the denominator is positive and independent of the action. Thus an individual leaf ranks actions by their aggregate measured training cost. The ensemble averages its 32 tree predictions. This construction does not establish a minimum cost on unseen requests.

For the fixed action $a_0$ selected for the request's propagator structure, an alternative $a_*$ is proposed only if
\begin{equation}
 \widehat y_{a_*}<(1-\eta)\widehat y_{a_0}.
 \label{eq:reconstruction-margin}
\end{equation}
The margin $\eta$ is a defined acceptance parameter. Validation compares $\eta=0,0.05,0.10$, tree depths 2, 4, and 6, and minimum leaf populations 2 and 4. The complete selected policy must cost less than 0.97 times its structure-dependent fixed control and select no failed validation action. The latest policy does not pass this condition. The final application therefore uses full reconstruction with eight evaluation points grouped together for the planar box, nonplanar box, and sunrise structures, and the preceding procedure for the three-loop vacuum structure.

To specify the features, define target excess power $D=\sum_j\max(a_j-1,0)$, numerator rank $R=\sum_j\max(-a_j,0)$, and line count $P_T=\sum_j{\bf1}_{a_j>0}$. A line is incident on a loop momentum when that momentum occurs in its inverse propagator. The 29 features are listed in Table~\ref{tab:reconstruction-features}. Every feature follows from the supplied indices and propagator definitions.
\begin{table}[!htbp]\centering\small
\caption{Inputs to the algebra-cost model. Means and extrema over targets use the complete requested batch.}
\label{tab:reconstruction-features}
\begin{tabularx}{\linewidth}{@{}lY@{}}\toprule
Positions & Quantities\\\midrule
1--6 & Number of loops, external momenta, integral coordinates, physical propagators, symbolic external invariants, and $\log(1+n_T)$ for $n_T$ targets\\
7--13 & Maximum and mean $D$, maximum and mean $R$, minimum and maximum $P_T$, and number of distinct sectors\\
14--18 & Number of massive coordinates; number of physical lines containing several loops; maximum loop incidence per physical line; maximum physical-line incidence per loop; mean external-momentum incidence per coordinate\\
19--22 & Minimum $D$, minimum $R$, mean $P_T$, and maximum absolute index\\
23--25 & Mean excess power on massive coordinates; mean numerator rank weighted by loop incidence; mean excess power weighted by loop incidence\\
26--29 & Maximum sum of positive powers; maximum number of absent physical lines; maximum numerator and denominator bit lengths among rational masses and scalar products\\\bottomrule
\end{tabularx}\end{table}
The ordering within positions 7--10 is maximum $D$, maximum $R$, mean $D$, mean $R$. The remaining entries follow the listed order. The bit lengths describe the supplied rational numbers and do not encode a family identifier.

\section{Physical test inputs}
\label{app:families-new}
All profiles below define squared quantities in a common fixed unit. The zero-invariant families use shifts $(0,p,q,r)$ and masses $(x,y,0,0)$, with
\begin{align}
 &p^2=0,\quad p\cdot q=-h-1,\quad p\cdot r=-h-2,\nonumber\\
 &q^2=-2h-7,\quad q\cdot r=-h-3,\quad r^2=-3h-13.
 \label{eq:light-input}
\end{align}
The unequal case has $(x,y)=(h+3,h+11)$, the equal case $(h+3,h+3)$, the mixed case $(h+3,0)$, and the massless case $(0,0)$. Targets have the form $(1+\delta_1,1+\delta_2,-b_1,-b_2)$, where $\delta_1+\delta_2$ is the excess positive power and $b_1+b_2$ is numerator rank.

Ordinary tensor families have $P$ physical lines, one auxiliary coordinate, shifts $q_1=0$, $q_{i+1}=p_i$, and
\begin{equation}
 m_i^2=h+11i^2\ (i\leq P),\quad m_{P+1}^2=0,\quad
 p_i^2=-(10h+17i),\quad p_i\cdot p_j=-(h+i+2j-2)\ (i<j).
 \label{eq:ordinary-input}
\end{equation}
Appendix~\ref{app:direct-inputs} specifies the current profiles, targets, and separate vacuum grids.

For the native planar box, let $p_1,p_2,p_3$ be massless, $2p_1\cdot p_2=s$, $2p_1\cdot p_3=t$, and $2p_2\cdot p_3=-s-t$. The seven physical line momenta are
\begin{equation}
 (\ell_1,\ell_1-p_1,\ell_1-p_1-p_2,\ell_2,
 \ell_2-p_1-p_2,\ell_2-p_1-p_2-p_3,\ell_1-\ell_2),
 \label{eq:planar-input}
\end{equation}
and the auxiliary momenta are $(\ell_1-p_3,\ell_2-p_1)$. The nonplanar routing follows Ref.~\cite{KiraOriginal}:
\begin{align}
 (Q_1,\ldots,Q_9)=(&\ell_1,\ell_2,p_1-\ell_1,p_2-\ell_2,
 p_1+p_2-p_3-\ell_1,\nonumber\\
 &p_1-\ell_1+\ell_2,p_1-p_3-\ell_1+\ell_2,
 \ell_1-p_3,\ell_2-p_1-p_2).
 \label{eq:nonplanar-input}
\end{align}
The first seven entries are physical and all masses vanish. Its convention gives $p_i^2=(p_1+p_2-p_3)^2=0$, $s=(p_1+p_2)^2$, and $t=(p_1-p_3)^2$. Expanding the squares gives $2p_1\cdot p_3=-t$ and $2p_2\cdot p_3=s+t$. The low-rank profiles use $(s,t)=(-1009-12h,-1021-14h)$ for the planar graph and $(-113-18h,-127-10h)$ for the nonplanar graph. Appendix~\ref{app:reconstruction-inputs} gives the independently chosen high-rank invariants.

Every high-rank box target is $(1,1,1,1,1,1,1,-u,-v)$ with $u,v\geq0$ and $1\leq u+v\leq4$. The current split uses the profiles and batch counts in Appendix~\ref{app:reconstruction-inputs}. Both graph types occur in all splits, so the test measures new batches and kinematic points within these structures.

\section{Time ratios and uncertainty}
\label{app:statistics}
For the final shared-call comparison, let $U_{mr}$ be elapsed time for the entire 93-request collection, method $m$, and repetition $r$. We define
\begin{equation}
 \widetilde U_m=\operatorname{median}_{r=1,2,3}U_{mr},\qquad
 \rho_{\rm collection}=\frac{\widetilde U_{\rm final}}{\widetilde U_{\rm preceding}},\qquad
 g=1-\rho_{\rm collection}.
 \label{eq:collection-ratio}
\end{equation}
The saving $g$ follows by dividing the time difference $\widetilde U_{\rm preceding}-\widetilde U_{\rm final}$ by the preceding time. The observed range is $[\min_r U_{mr},\max_r U_{mr}]$. It describes the three measurements and is not a statistical confidence interval. Internal per-request timers exclude some initialization and collection overhead, so summing their medians need not give $\widetilde U_m$.

Let $T_{imr}$ be complete elapsed time for request $i$, method $m$, and repetition $r$. For the set $\mathcal C$ completed by both methods in all repetitions, define
\begin{equation}
 S_m=\sum_{i\in\mathcal C}\operatorname{median}_{r=1,2,3}T_{imr},\qquad
 \rho=\frac{S_{\rm new}}{S_{\rm control}},\qquad
 \text{fractional saving}=1-\rho.
 \label{eq:time-ratio}
\end{equation}
These definitions compare identical completed requests. Dividing both times by the control time gives $\rho$, and subtracting gives the fractional saving. Completion counts remain separate because an unfinished reduction has no verified finite speedup. To estimate variation across requests, we draw $|\mathcal C|$ paired request indices with replacement and recompute $\rho$. We report the central 95\% of 5,000 such ratios, using the bootstrap construction of Ref.~\cite{Efron1979}. This interval describes the sampled request collection and does not establish performance for every integral family.


For the historical compiled-row result, the estimator is the geometric mean of paired ratios, $\rho_G=\exp[n^{-1}\sum_{i=1}^{n}\log(t_i^{\rm learned}/t_i^{\rm native})]$, where $n$ is the number of matched units in that study. This is the $n$th root of their product, obtained by applying the logarithm to that product. It weights relative changes equally, unlike the total-time ratio in Eq.~\eqref{eq:time-ratio}; the notes preserve the original grouping and resampling procedure.
